% Options for packages loaded elsewhere
\PassOptionsToPackage{unicode}{hyperref}
\PassOptionsToPackage{hyphens}{url}
\documentclass[
  ignorenonframetext,
]{beamer}
\newif\ifbibliography
\usepackage{pgfpages}
\setbeamertemplate{caption}[numbered]
\setbeamertemplate{caption label separator}{: }
\setbeamercolor{caption name}{fg=normal text.fg}
\beamertemplatenavigationsymbolsempty
% remove section numbering
\setbeamertemplate{part page}{
  \centering
  \begin{beamercolorbox}[sep=16pt,center]{part title}
    \usebeamerfont{part title}\insertpart\par
  \end{beamercolorbox}
}
\setbeamertemplate{section page}{
  \centering
  \begin{beamercolorbox}[sep=12pt,center]{section title}
    \usebeamerfont{section title}\insertsection\par
  \end{beamercolorbox}
}
\setbeamertemplate{subsection page}{
  \centering
  \begin{beamercolorbox}[sep=8pt,center]{subsection title}
    \usebeamerfont{subsection title}\insertsubsection\par
  \end{beamercolorbox}
}
% Prevent slide breaks in the middle of a paragraph
\widowpenalties 1 10000
\raggedbottom
\AtBeginPart{
  \frame{\partpage}
}
\AtBeginSection{
  \ifbibliography
  \else
    \frame{\sectionpage}
  \fi
}
\AtBeginSubsection{
  \frame{\subsectionpage}
}
\usepackage{iftex}
\ifPDFTeX
  \usepackage[T1]{fontenc}
  \usepackage[utf8]{inputenc}
  \usepackage{textcomp} % provide euro and other symbols
\else % if luatex or xetex
  \usepackage{unicode-math} % this also loads fontspec
  \defaultfontfeatures{Scale=MatchLowercase}
  \defaultfontfeatures[\rmfamily]{Ligatures=TeX,Scale=1}
\fi
\usepackage{lmodern}
\ifPDFTeX\else
  % xetex/luatex font selection
\fi
% Use upquote if available, for straight quotes in verbatim environments
\IfFileExists{upquote.sty}{\usepackage{upquote}}{}
\IfFileExists{microtype.sty}{% use microtype if available
  \usepackage[]{microtype}
  \UseMicrotypeSet[protrusion]{basicmath} % disable protrusion for tt fonts
}{}
\makeatletter
\@ifundefined{KOMAClassName}{% if non-KOMA class
  \IfFileExists{parskip.sty}{%
    \usepackage{parskip}
  }{% else
    \setlength{\parindent}{0pt}
    \setlength{\parskip}{6pt plus 2pt minus 1pt}}
}{% if KOMA class
  \KOMAoptions{parskip=half}}
\makeatother
\usepackage{longtable,booktabs,array}
\newcounter{none} % for unnumbered tables
\usepackage{calc} % for calculating minipage widths
\usepackage{caption}
% Make caption package work with longtable
\makeatletter
\def\fnum@table{\tablename~\thetable}
\makeatother
\setlength{\emergencystretch}{3em} % prevent overfull lines
\providecommand{\tightlist}{%
  \setlength{\itemsep}{0pt}\setlength{\parskip}{0pt}}
\usepackage{bookmark}
\IfFileExists{xurl.sty}{\usepackage{xurl}}{} % add URL line breaks if available
\urlstyle{same}
\hypersetup{
  hidelinks,
  pdfcreator={LaTeX via pandoc}}

\author{\texorpdfstring{}{}}
\date{}

\begin{document}

\begin{frame}[fragile]{Infrastructure Module Reference}
\protect\phantomsection\label{infrastructure-module-reference}
This section provides a detailed reference for all 13 infrastructure
subpackages, documenting their purpose, key classes, public API, and
integration points within the pipeline. The infrastructure layer
comprises \textasciitilde277 Python modules validated by 6,028 tests.
Each subpackage follows the Documentation Duality standard: every module
directory contains both an \texttt{AGENTS.md} machine-readable
specification and a \texttt{README.md} human-readable guide.

\begin{block}{\texttt{infrastructure.core} (71 modules)}
\protect\phantomsection\label{infrastructure.core-71-modules}
\textbf{Purpose}: Foundation utilities providing the bedrock services
consumed by all other modules and all projects.

\textbf{Key Components}:

{\def\LTcaptype{none} % do not increment counter
\begin{longtable}[]{@{}
  >{\raggedright\arraybackslash}p{(\linewidth - 2\tabcolsep) * \real{0.5500}}
  >{\raggedright\arraybackslash}p{(\linewidth - 2\tabcolsep) * \real{0.4500}}@{}}
\toprule\noalign{}
\begin{minipage}[b]{\linewidth}\raggedright
Component
\end{minipage} & \begin{minipage}[b]{\linewidth}\raggedright
Purpose
\end{minipage} \\
\midrule\noalign{}
\endhead
\texttt{logging\_utils.py} & Structured logger factory
(\texttt{get\_logger}) with colored console output and file rotation \\
\texttt{config\_loader.py} & YAML config parser (\texttt{load\_config})
with schema validation and default merging \\
\texttt{exceptions.py} & Exception hierarchy: \texttt{TemplateError} →
\texttt{ConfigurationError}, \texttt{ValidationError},
\texttt{BuildError} \\
\texttt{environment.py} & Environment detection, Python command
resolution, \texttt{PYTHONPATH} management \\
\texttt{progress.py} & \texttt{ProgressBar} for pipeline stage progress
reporting \\
\texttt{checkpoint.py} & \texttt{CheckpointManager} for resumable
pipeline execution \\
\texttt{health.py} & \texttt{SystemHealthChecker} for pre-pipeline
dependency validation \\
\texttt{performance.py} & \texttt{monitor\_performance} context manager
for timing and memory tracking \\
\texttt{\_optional\_deps.py} & Lazy loading of optional dependencies
(\texttt{psutil}, \texttt{reportlab}) \\
\bottomrule\noalign{}
\end{longtable}
}

\textbf{Integration}: Every module and script imports from
\texttt{core}. The exception hierarchy is used for pipeline flow
control---\texttt{ValidationError} triggers stage failure,
\texttt{BuildError} halts the entire pipeline. The lazy loader in
\texttt{\_optional\_deps.py} separates core imports from optional
subpackages, preventing cascading import failures in environments with
partial dependency sets.
\end{block}

\begin{block}{\texttt{infrastructure.documentation} (6 modules)}
\protect\phantomsection\label{infrastructure.documentation-6-modules}
\textbf{Purpose}: Documentation management, figure registration, and API
glossary generation.

\textbf{Key Components}:

{\def\LTcaptype{none} % do not increment counter
\begin{longtable}[]{@{}
  >{\raggedright\arraybackslash}p{(\linewidth - 2\tabcolsep) * \real{0.5500}}
  >{\raggedright\arraybackslash}p{(\linewidth - 2\tabcolsep) * \real{0.4500}}@{}}
\toprule\noalign{}
\begin{minipage}[b]{\linewidth}\raggedright
Component
\end{minipage} & \begin{minipage}[b]{\linewidth}\raggedright
Purpose
\end{minipage} \\
\midrule\noalign{}
\endhead
\texttt{figure\_manager.py} & \texttt{FigureManager} --- maintains a
JSON registry of all generated figures with captions, labels, and
generation metadata \\
\texttt{glossary\_gen.py} & \texttt{generate\_glossary} --- programmatic
API glossary extraction from Python source files \\
\bottomrule\noalign{}
\end{longtable}
}

\textbf{Integration}: Called by project scripts during Stage 02 to
register figures for automated cross-referencing in the manuscript. The
glossary generator supports the Documentation Duality standard by
extracting docstrings and function signatures.
\end{block}

\begin{block}{\texttt{infrastructure.llm} (53 modules)}
\protect\phantomsection\label{infrastructure.llm-53-modules}
\textbf{Purpose}: Local LLM integration for automated manuscript review,
translation, and literature search.

\textbf{Key Components}:

{\def\LTcaptype{none} % do not increment counter
\begin{longtable}[]{@{}
  >{\raggedright\arraybackslash}p{(\linewidth - 2\tabcolsep) * \real{0.5500}}
  >{\raggedright\arraybackslash}p{(\linewidth - 2\tabcolsep) * \real{0.4500}}@{}}
\toprule\noalign{}
\begin{minipage}[b]{\linewidth}\raggedright
Component
\end{minipage} & \begin{minipage}[b]{\linewidth}\raggedright
Purpose
\end{minipage} \\
\midrule\noalign{}
\endhead
\texttt{review.py} & Executive summary and quality review generation \\
\texttt{translation.py} & Abstract translation into configured target
languages \\
\texttt{client.py} & Ollama HTTP client with retry logic and timeout
management \\
\texttt{literature/} & Literature search subpackage with semantic query
support \\
\texttt{templates/} & Prompt templates for structured LLM
interactions \\
\bottomrule\noalign{}
\end{longtable}
}

\textbf{Integration}: Invoked during Stage 06. Requires a running Ollama
instance. Gracefully degrades when unavailable. The literature search
subpackage enables programmatic discovery of related work during
manuscript preparation.
\end{block}

\begin{block}{\texttt{infrastructure.project} (5 modules)}
\protect\phantomsection\label{infrastructure.project-5-modules}
\textbf{Purpose}: Project discovery and workspace management.

\textbf{Key Components}:

{\def\LTcaptype{none} % do not increment counter
\begin{longtable}[]{@{}
  >{\raggedright\arraybackslash}p{(\linewidth - 2\tabcolsep) * \real{0.5500}}
  >{\raggedright\arraybackslash}p{(\linewidth - 2\tabcolsep) * \real{0.4500}}@{}}
\toprule\noalign{}
\begin{minipage}[b]{\linewidth}\raggedright
Component
\end{minipage} & \begin{minipage}[b]{\linewidth}\raggedright
Purpose
\end{minipage} \\
\midrule\noalign{}
\endhead
\texttt{discovery.py} & \texttt{\_discover\_project} --- finds valid
project directories by scanning for \texttt{manuscript/config.yaml} \\
\texttt{workspace.py} & Workspace initialization and cleanup
utilities \\
\bottomrule\noalign{}
\end{longtable}
}

\textbf{Integration}: Used by \texttt{execute\_pipeline.py} and
\texttt{run.sh} to identify which projects can be built. The discovery
algorithm enforces the Standalone Project Paradigm: a directory is a
valid project if and only if it contains
\texttt{manuscript/config.yaml}.
\end{block}

\begin{block}{\texttt{infrastructure.publishing} (14 modules)}
\protect\phantomsection\label{infrastructure.publishing-14-modules}
\textbf{Purpose}: Academic publishing metadata and citation generation.

\textbf{Key Components}:

{\def\LTcaptype{none} % do not increment counter
\begin{longtable}[]{@{}
  >{\raggedright\arraybackslash}p{(\linewidth - 2\tabcolsep) * \real{0.5500}}
  >{\raggedright\arraybackslash}p{(\linewidth - 2\tabcolsep) * \real{0.4500}}@{}}
\toprule\noalign{}
\begin{minipage}[b]{\linewidth}\raggedright
Component
\end{minipage} & \begin{minipage}[b]{\linewidth}\raggedright
Purpose
\end{minipage} \\
\midrule\noalign{}
\endhead
\texttt{models.py} & \texttt{PublicationMetadata} dataclass with direct
attribute access and dynamic \texttt{getattr} fallback \\
\texttt{citations.py} & \texttt{generate\_citation\_apa},
\texttt{generate\_citation\_bibtex}, \texttt{generate\_citation\_mla} \\
\texttt{zenodo.py} & Zenodo API integration for DOI registration \\
\bottomrule\noalign{}
\end{longtable}
}

\textbf{Integration}: Used during Stage 02 by analysis scripts to
extract publishable metadata from results. Citation generators produce
correctly formatted strings from \texttt{config.yaml} metadata.
\end{block}

\begin{block}{\texttt{infrastructure.rendering} (27 modules)}
\protect\phantomsection\label{infrastructure.rendering-27-modules}
\textbf{Purpose}: Multi-format document rendering (Markdown → LaTeX →
PDF, HTML reports).

\textbf{Key Components}:

{\def\LTcaptype{none} % do not increment counter
\begin{longtable}[]{@{}
  >{\raggedright\arraybackslash}p{(\linewidth - 2\tabcolsep) * \real{0.5500}}
  >{\raggedright\arraybackslash}p{(\linewidth - 2\tabcolsep) * \real{0.4500}}@{}}
\toprule\noalign{}
\begin{minipage}[b]{\linewidth}\raggedright
Component
\end{minipage} & \begin{minipage}[b]{\linewidth}\raggedright
Purpose
\end{minipage} \\
\midrule\noalign{}
\endhead
\texttt{pandoc.py} & Pandoc invocation with custom filters and metadata
injection \\
\texttt{latex.py} & XeLaTeX compilation with auxiliary file management
and stale \texttt{.aux} cleanup \\
\texttt{pdf\_builder.py} & End-to-end PDF construction orchestrating
Pandoc and XeLaTeX \\
\texttt{html\_report.py} & HTML executive report generation \\
\texttt{markdown\_report.py} & Markdown-format report generation \\
\bottomrule\noalign{}
\end{longtable}
}

\textbf{Integration}: Core of Stage 03. Reads \texttt{manuscript/*.md}
and \texttt{config.yaml}, produces
\texttt{output/\textless{}project\textgreater{}.pdf}. The auxiliary file
cleanup resolves a known rendering hazard where stale \texttt{.aux}
files cause ``Division by 0'' LaTeX errors.
\end{block}

\begin{block}{\texttt{infrastructure.reporting} (45 modules)}
\protect\phantomsection\label{infrastructure.reporting-45-modules}
\textbf{Purpose}: Pipeline reporting, test result aggregation, and
coverage analysis.

\textbf{Key Components}:

{\def\LTcaptype{none} % do not increment counter
\begin{longtable}[]{@{}
  >{\raggedright\arraybackslash}p{(\linewidth - 2\tabcolsep) * \real{0.5500}}
  >{\raggedright\arraybackslash}p{(\linewidth - 2\tabcolsep) * \real{0.4500}}@{}}
\toprule\noalign{}
\begin{minipage}[b]{\linewidth}\raggedright
Component
\end{minipage} & \begin{minipage}[b]{\linewidth}\raggedright
Purpose
\end{minipage} \\
\midrule\noalign{}
\endhead
\texttt{coverage\_parser.py} & Cascading parse strategies for pytest
output: \texttt{\_parse\_failures\_section},
\texttt{\_parse\_failures\_verbose}, \texttt{\_parse\_failures\_short},
\texttt{\_parse\_failures\_timeout},
\texttt{\_parse\_failures\_fallback} \\
\texttt{report\_generator.py} & Executive report generation in JSON and
Markdown \\
\texttt{statistics.py} & \texttt{collect\_output\_statistics} ---
enumerates output directory contents \\
\bottomrule\noalign{}
\end{longtable}
}

\textbf{Integration}: Used during Stages 01, 04, and 07 for test result
parsing, validation reporting, and executive summary generation. The
cascading parser handles all pytest output formats robustly, falling
through five strategies to ensure no test failure is silently missed.
\end{block}

\begin{block}{\texttt{infrastructure.scientific} (6 modules)}
\protect\phantomsection\label{infrastructure.scientific-6-modules}
\textbf{Purpose}: Scientific computing utilities for numerical analysis
and benchmarking.

\textbf{Key Components}:

{\def\LTcaptype{none} % do not increment counter
\begin{longtable}[]{@{}
  >{\raggedright\arraybackslash}p{(\linewidth - 2\tabcolsep) * \real{0.5500}}
  >{\raggedright\arraybackslash}p{(\linewidth - 2\tabcolsep) * \real{0.4500}}@{}}
\toprule\noalign{}
\begin{minipage}[b]{\linewidth}\raggedright
Component
\end{minipage} & \begin{minipage}[b]{\linewidth}\raggedright
Purpose
\end{minipage} \\
\midrule\noalign{}
\endhead
\texttt{stability.py} & \texttt{check\_numerical\_stability} --- tests
functions for NaN/Inf behavior across input ranges \\
\texttt{benchmarking.py} & \texttt{benchmark\_function} --- measures
execution time and memory usage \\
\texttt{simulation.py} & Scientific simulation framework with parameter
sweeps \\
\bottomrule\noalign{}
\end{longtable}
}

\textbf{Integration}: Used by \texttt{code\_project}'s analysis scripts
during Stage 02 for algorithm validation. The stability checker is
critical for ensuring that numerical results are reproducible across
different floating-point environments.
\end{block}

\begin{block}{\texttt{infrastructure.steganography} (10 modules)}
\protect\phantomsection\label{infrastructure.steganography-10-modules}
\textbf{Purpose}: Cryptographic watermarking and provenance embedding
for PDF artifacts.

\textbf{Key Components}:

{\def\LTcaptype{none} % do not increment counter
\begin{longtable}[]{@{}
  >{\raggedright\arraybackslash}p{(\linewidth - 2\tabcolsep) * \real{0.5500}}
  >{\raggedright\arraybackslash}p{(\linewidth - 2\tabcolsep) * \real{0.4500}}@{}}
\toprule\noalign{}
\begin{minipage}[b]{\linewidth}\raggedright
Component
\end{minipage} & \begin{minipage}[b]{\linewidth}\raggedright
Purpose
\end{minipage} \\
\midrule\noalign{}
\endhead
\texttt{metadata.py} & \texttt{inject\_pdf\_metadata} --- embeds XMP
metadata and PDF Info dictionary entries \\
\texttt{config.py} & \texttt{DocumentMetadata} dataclass for
steganography configuration \\
\texttt{overlay.py} & Alpha-channel text overlay with build timestamp
and commit hash \\
\texttt{qr.py} & QR code generation and injection into PDF pages \\
\texttt{hash.py} & SHA-256/SHA-512 hash computation for tamper
detection \\
\bottomrule\noalign{}
\end{longtable}
}

\textbf{Integration}: Invoked by \texttt{secure\_run.sh} after the main
pipeline completes. Reads the rendered PDF and produces a
steganographically watermarked copy with an accompanying
\texttt{.hashes.json} manifest.
\end{block}

\begin{block}{\texttt{infrastructure.validation} (36 modules)}
\protect\phantomsection\label{infrastructure.validation-36-modules}
\textbf{Purpose}: Quality assurance and integrity verification for all
pipeline artifacts.

\textbf{Key Components}:

{\def\LTcaptype{none} % do not increment counter
\begin{longtable}[]{@{}
  >{\raggedright\arraybackslash}p{(\linewidth - 2\tabcolsep) * \real{0.5500}}
  >{\raggedright\arraybackslash}p{(\linewidth - 2\tabcolsep) * \real{0.4500}}@{}}
\toprule\noalign{}
\begin{minipage}[b]{\linewidth}\raggedright
Component
\end{minipage} & \begin{minipage}[b]{\linewidth}\raggedright
Purpose
\end{minipage} \\
\midrule\noalign{}
\endhead
\texttt{pdf\_validator.py} & PDF structural integrity checking (xref
table, trailer, page count) \\
\texttt{markdown\_validator.py} & Markdown linting (heading hierarchy,
link integrity, orphan references) \\
\texttt{integrity.py} & \texttt{verify\_output\_integrity} ---
comprehensive output directory validation \\
\texttt{cli.py} & Command-line interface for standalone validation
operations \\
\bottomrule\noalign{}
\end{longtable}
}

\textbf{Integration}: Core of Stage 04. Validates all generated
artifacts before they are finalized. The validation module is the most
module-dense package (36 files), reflecting the breadth of integrity
checks required across PDF, Markdown, image, and manifest formats.
\end{block}

\begin{block}{\texttt{infrastructure.skills} (4 modules)}
\protect\phantomsection\label{infrastructure.skills-4-modules}
\textbf{Purpose}: Programmatic discovery and manifest generation for
AI-agent skill descriptors.

\textbf{Key Components}:

{\def\LTcaptype{none} % do not increment counter
\begin{longtable}[]{@{}
  >{\raggedright\arraybackslash}p{(\linewidth - 2\tabcolsep) * \real{0.5500}}
  >{\raggedright\arraybackslash}p{(\linewidth - 2\tabcolsep) * \real{0.4500}}@{}}
\toprule\noalign{}
\begin{minipage}[b]{\linewidth}\raggedright
Component
\end{minipage} & \begin{minipage}[b]{\linewidth}\raggedright
Purpose
\end{minipage} \\
\midrule\noalign{}
\endhead
\texttt{discovery.py} & \texttt{discover\_skills} --- scans
\texttt{SKILL.md} files across infrastructure modules and generates a
unified skill manifest \\
\texttt{cli.py} & Command-line interface for writing and validating
skill manifests \\
\bottomrule\noalign{}
\end{longtable}
}

\textbf{Integration}: Maintains \texttt{.cursor/skill\_manifest.json}
for IDE-level agent discovery. The skill manifest aggregates all
\texttt{SKILL.md} frontmatter into a single JSON index, enabling editors
and AI coding assistants to locate module capabilities without
traversing the directory tree.
\end{block}

\begin{block}{Infrastructure Maturity Summary}
\protect\phantomsection\label{infrastructure-maturity-summary}
The 13-module architecture achieves 100\% Tier 1--2 documentation
coverage (\texttt{AGENTS.md}, \texttt{README.md}) with Tier-3
\texttt{SKILL.md} skill descriptors across all 10 active subpackages,
83\%+ aggregate test coverage (exceeding the 60\% infrastructure
threshold by a wide margin), and zero mock-object violations. Every
active module exposes a machine-readable skill descriptor aligned with
the Model Context Protocol {[}@anthropic2024mcp{]}, making the
infrastructure layer not merely documented but \emph{programmatically
discoverable}---a prerequisite for the agentic research automation
paradigm described in the
\href{03c_documentation.md\#documentation-duality-and-ai-collaboration}{Documentation
Duality} and
\href{05d_ai_collaboration.md\#the-ai-collaboration-model}{AI
Collaboration} sections. The combination of high coverage, complete
documentation, and protocol-aligned discoverability positions
\texttt{template/}'s infrastructure as deployment-ready research
software rather than a prototype, satisfying the executability and
metadata quality indicators defined by Garijo et al.'s FAIRsoft
evaluator {[}@garijo2024fairsoft{]}.
\end{block}
\end{frame}

\end{document}
